\chapter{Software Requirements Specification (SRS)}

\section{Purpose and Scope of the SRS}
The purpose of this Software Requirements Specification (SRS) is to establish a rigorous, comprehensive, and unambiguous technical specification for the \textbf{GrievanceIQ} platform. This document delineates both the functional capabilities and the non-functional constraints governing system design, architecture, interface integration, security compliance, and operational deployment. It serves as the authoritative engineering benchmark for system verification, validation, and academic appraisal.

\section{Intended Users and Stakeholders}
\begin{enumerate}
    \item \textbf{Bank Customers / Aggrieved Depositors:} Individuals filing disputes, tracking case progression, dictating voice complaints, and downloading statutory receipts.
    \item \textbf{Grievance Redressal Officers (Resolvers):} Bank personnel evaluating evidence documents, verifying claim veracity, reviewing AI decision traces, and executing formal resolution memos.
    \item \textbf{Branch Managers \& Administrative Executives:} Supervisory authorities monitoring SLA adherence, inspecting officer productivity, simulating capacity surges, and auditing DPDP Act compliance.
    \item \textbf{Regulatory Authorities (e.g., RBI Ombudsman, Auditors):} Independent oversight bodies auditing resolution fairness, average time-to-redress, and consumer protection adherence.
\end{enumerate}

\section{Overall Product Description}
GrievanceIQ is an enterprise-grade, decoupled three-tier software architecture. The frontend is built as a single-page reactive application using React 18, Vite, and Tailwind CSS. The backend is an asynchronous RESTful microservice powered by Python FastAPI, Uvicorn (ASGI), and SQLite. The platform provides real-time speech processing via browser-native Web Speech interfaces, dynamic 2D barcode generation via ReportLab, and machine-learning text classification via Scikit-Learn.

\section{Functional Requirements}
The core operational requirements of GrievanceIQ are formalized in Table~\ref{tab:functional_requirements}.

\begin{table}[p]
\centering
\caption{Formal Functional Requirements Specification (FR-01 to FR-10)}
\label{tab:functional_requirements}
\begin{tabular}{|p{1.4cm}|p{3.2cm}|p{3.2cm}|p{3.2cm}|p{3.4cm}|}
\hline
\textbf{Req. ID} & \textbf{Requirement Title} & \textbf{Input} & \textbf{Expected Processing} & \textbf{Output / Acceptance Criteria} \\ \hline
\textbf{FR-01} & Trilingual Complaint Registration & Customer details, text description, bank, transaction ID & Ingestion parsing, trilingual encoding check & Generated Complaint ID (e.g., GIQ-2026-X) within 200\,ms \\ \hline
\textbf{FR-02} & Autonomous ML Categorization & Plain-text complaint description & TF-IDF vectorization + Multinomial Naive Bayes & Category assignment with confidence score $\ge 0.50$ \\ \hline
\textbf{FR-03} & Real-Time Voice Dictation & Spoken audio stream (Microphone) & W3C Web Speech API event handling & Real-time transcript with zero duplicate lines \\ \hline
\textbf{FR-04} & DPDP Act 2023 PII Masking & Raw text containing Card/Aadhaar/PAN & Regex token replacement (\texttt{XXXX-XXXX-...}) & Scrubbed text; zero unmasked PII saved to DB \\ \hline
\textbf{FR-05} & Statutory PDF Receipt Generation & Complaint metadata, grievance ID & ReportLab canvas layout, dynamic QR encoding & Valid PDF download ($< 150\,\text{KB}$) with scannable QR \\ \hline
\textbf{FR-06} & Document Evidence OCR Verification & Uploaded image (PNG/JPG) or PDF bill & OCR text extraction + 4-point entity matching & Verification score (0--100\%) and entity match badges \\ \hline
\textbf{FR-07} & Predictive Queue Modeling ($M/M/c$) & Arrival rate $\lambda$, service rate $\mu$, servers $c$ & Analytical Erlang-C queueing computation & Dwell time $W_q$, utilization $\rho$, and breach probability \\ \hline
\textbf{FR-08} & Dynamic SLA Tracking \& Alerts & Complaint timestamps, priority tier & Time delta comparison against 72h / 24h limits & Visual color-coded countdowns and breach badges \\ \hline
\textbf{FR-09} & AI Batch \& Manual Redressal & Resolution memo text, officer signature & State transition from Pending $\rightarrow$ Resolved & Audit log entry, resolution timestamp, memo saved \\ \hline
\textbf{FR-10} & Trilingual Dynamic UI Localization & Selected language state (EN/HI/MR) & React Context provider key lookup & Instantaneous UI language swap across 222 keys \\ \hline
\end{tabular}
\end{table}

\section{Non-Functional Requirements}
\begin{itemize}
    \item \textbf{Performance:} All transactional REST endpoints must respond within 150\,ms under standard operational conditions. Analytical computation endpoints must respond within 50\,ms through memory-optimized aggregations.
    \item \textbf{Scalability:} The architectural data layer and async event loops must sustain at least 100 concurrent HTTP requests without connection drops or resource exhaustion.
    \item \textbf{Availability and Reliability:} The backend must maintain 99.9\% uptime. Transactional database operations must adhere strictly to ACID properties, utilizing explicit commit/rollback mechanisms.
    \item \textbf{Security and Privacy:}
    \begin{itemize}
        \item Full compliance with India's DPDP Act 2023: customer Aadhaar numbers, 16-digit debit/credit card numbers, and Permanent Account Numbers must be masked prior to persistence.
        \item Role-Based Access Control (RBAC) ensuring customer complaints can only be resolved by authorized bank officers.
        \item Zero persistent plain-text logging of sensitive customer credentials.
    \end{itemize}
    \item \textbf{Usability and Accessibility:} The client interface must support high-contrast dark and light visual modes, conform to trilingual standards across English, Hindi, and Marathi, and provide accessible keyboard navigation.
    \item \textbf{Maintainability and Portability:} Decoupled frontend and backend architectures packaged with standardized configuration files (\texttt{vite.config.js}, \texttt{requirements.txt}), enabling rapid deployment across Windows, Linux, and containerized Docker environments.
\end{itemize}

\section{Hardware and Software Requirements}
\subsection{Hardware Requirements}
\begin{itemize}
    \item \textbf{Processor:} Quad-Core 64-bit x86/ARM Processor (Intel Core i5 8th Gen / AMD Ryzen 5 or equivalent).
    \item \textbf{Memory (RAM):} Minimum 8.0 GB (16.0 GB recommended for concurrent OCR model execution).
    \item \textbf{Storage:} Minimum 2.0 GB free solid-state storage (SSD) for database records, models, and virtual environments.
    \item \textbf{Audio Input:} Standard microphone hardware supporting 16 kHz or 44.1 kHz PCM audio capture for voice dictation.
\end{itemize}

\subsection{Software Requirements}
\begin{itemize}
    \item \textbf{Operating System:} Microsoft Windows 10/11, Ubuntu Linux 22.04 LTS, or macOS 13+.
    \item \textbf{Backend Runtime:} Python 3.10.x or 3.11.x 64-bit.
    \item \textbf{Web Framework:} FastAPI v0.100+ with Uvicorn ASGI Web Server.
    \item \textbf{Frontend Runtime:} Node.js v18.x or v20.x LTS, with NPM v9.x+.
    \item \textbf{Frontend Framework:} React 18.2, Vite 5.x, Tailwind CSS 3.x, Lucide React Icons.
    \item \textbf{Machine Learning \& Data Science:} Scikit-Learn, NumPy, Pandas, Joblib.
    \item \textbf{Document \& Graphics Processing:} ReportLab 4.x, Pillow 10.x, EasyOCR.
    \item \textbf{Browser Compatibility:} Google Chrome 100+, Microsoft Edge 100+ (for full W3C Web Speech API support).
\end{itemize}

\section{Design and Implementation Constraints}
\begin{itemize}
    \item The speech recognition pipeline is governed by browser permission models and requires active user microphone authorization.
    \item ReportLab PDF canvas coordinates require deterministic typography calculations to guarantee zero overflow across multi-page complaint descriptions.
    \item SQLite database concurrent write concurrency is managed via WAL (Write-Ahead Logging) mode to prevent table locking anomalies during simultaneous grievance submissions.
\end{itemize}
